\begin{abstract}

%% ==================== 通用摘要模板 (v1.5.5 重写, 通用化) ====================
%%
%% 设计原则:
%% 1. 整段都是 \TeX 注释 (以 % 开头), check 15 跳过 % 行 → dryrun 不报 RED
%% 2. 留 1 行 INLINE 占位符 【】, check 15 抓【】→ 跑题用户填完才 GREEN
%% 3. 跑题用户复制本文件后, 删注释 + 填摘要正文, 删 【】占位符行
%%
%% 字数: 800-1000 字, 1 页 A4
%% 结构 (v1.5.3 板块自查 01 借鉴 BZD 7 维诊断):
%%   - 第 1 段 (2 句): 背景 + 核心问题 (不堆砌数据/算法)
%%   - 主体 N 段: 一问一段, 每段含 [任务 → 数据/假设 → 模型+数学类型 → 算法 → 量化结果 → 结论]
%%   - 最后段: 创新点 + 检验方式 + 优势 + 推广价值 + 局限 + 改进方向
%%   - 关键词: 4-6 个, 3 方向 (问题/模型/算法创新)
%% 写法: 第三人称, 不写图表/公式/参考文献编号, 数值必须来自正文不得虚构

%%\textbf{第一段 (引言 2 句):} 随着【你的研究背景, 1 句】, 【你的研究对象/核心问题, 1 句】。本文基于【数据概况, 1082 例】开展【研究目标, 4 问研究】。

%%\textbf{针对问题一 (按你实际题改):} 为【任务, 1 句】, 采用【模型名 + 数学类型, 如: 多元线性回归 (MLR) + 随机森林回归 (RFR)】。结果表明【关键特征 + 量化结果】。

%%\textbf{针对问题二 (按你实际题改):} 以【异常判据】为基准, 建立【模型】。推荐【最佳时点/参数】。

%%\textbf{针对问题三 (按你实际题改):} 按【严格按题面给的分界, 如: 题面 [20,28) [28,32) [32,36) [36,40) ≥40】分层, 采用【模型】。结果表明【关键结论】。

%%\textbf{针对问题四 (按你实际题改):} 对【数据, 严格按题面指定 sheet】, 构建【模型】。推荐【决策规则 + 比例】。

%%\textbf{创新点:} (1) 【模型层创新】; (2) 【方法层创新】。经【检验方式】, 模型有较好的稳健性; 局限: 【1 句】。

%%\keywords{【关键词 1, 问题方向】;【关键词 2, 模型方向】;【关键词 3, 算法方向】;【关键词 4, 创新点】;【关键词 5, 应用方向】}

\textbf{【在这里粘贴你的 800-1000 字摘要正文, 严格按上面 6 段结构 (引言 + 4 问 + 创新) + 关键词。删掉这一行, 替换为你的实际摘要。本行留【】= dryrun check 15 RED, 必须替换。】}

\label{abstract:end}

\end{abstract}
